\chapter{Em Có Đang Được Thương Không}
\markboth{Em Có Đang Được Thương Không}{Am I Still Being Loved}

\section{Điểm Kém Xong Em Nhìn Mặt Người Lớn}
\begin{truyen}{Điểm Kém Xong Em Nhìn Mặt Người Lớn}{After a Bad Grade, the Child Looks at Adult Faces}
\chuhoa{C}{ó những em nhận bài kiểm tra xong không nhìn lại lỗi trước. Các em nhìn mặt thầy cô, rồi nhìn mặt cha mẹ.}
\textit{(Some students do not first look at the mistakes after getting a test back. They look at the teacher's face, then at their parents' faces.)}

Điều các em sợ không chỉ là con điểm.
\textit{(What they fear is not only the score.)}

Có em hỏi rất nhỏ: “Con điểm này thì tối nay ở nhà có yên không thầy?”
\textit{(Some students ask very softly, “With this score, will things be peaceful at home tonight?”)}

Các em đang dò xem với kết quả này, mình còn đáng được dịu dàng như cũ không.
\textit{(They are checking whether they are still worthy of the same gentleness.)}

Có những đứa trẻ sống trong cảm giác tình thương lên xuống theo kết quả.
\textit{(Some children live with the feeling that love rises and falls with performance.)}

\end{truyen}

\begin{baihoc}
Một đứa trẻ cần biết rằng mình có thể sai mà không bị rút mất tình thương.
\textit{(A child needs to know they can fail without losing love.)}
\end{baihoc}

\begin{ghinhoanh}
\textbf{Short English to remember:}

A child needs to know they can fail without losing love.

\end{ghinhoanh}

\begin{tuhocnhanh}
\textbf{Gợi ý để dạy và tự nghĩ / Teaching and reflection prompts}

1. Vì sao trẻ nhìn mặt người lớn sau khi xem điểm? \textit{Why do children look at adult faces after seeing a grade?}

2. Điều trẻ đang sợ thật sự là gì? \textit{What are they truly afraid of?}

3. Mình đang gắn tình thương với kết quả ở mức nào? \textit{How much are we tying love to performance?}
\end{tuhocnhanh}
\ngancach

\section{Khi Em Ốm, Người Lớn Có Dịu Lại Không}
\begin{truyen}{Khi Em Ốm, Người Lớn Có Dịu Lại Không}{When I Am Unwell, Do Adults Become Gentle}
\chuhoa{C}{ó những em chỉ khi bệnh, khi ngất, khi khóc quá mới được người lớn nói nhỏ lại.}
\textit{(Some children only receive a softer tone from adults when they are sick, faint, or crying hard.)}

Các em sẽ học rất nhanh rằng phải yếu đi thật rõ thì mới được thương thấy được.
\textit{(They learn quickly that they must visibly weaken before love becomes visible.)}

Đó là một bài học buồn.
\textit{(That is a sad lesson.)}

Trẻ con không nên phải vỡ ra mới được đối xử nhẹ tay.
\textit{(Children should not have to break down before receiving gentleness.)}

\end{truyen}

\begin{baihoc}
Sự dịu dàng cần có trước khi một đứa trẻ kiệt sức, không phải chỉ sau đó.
\textit{(Gentleness should arrive before a child is exhausted, not only after.)}
\end{baihoc}

\begin{ghinhoanh}
\textbf{Short English to remember:}

Gentleness should arrive before a child breaks down.

\end{ghinhoanh}

\begin{tuhocnhanh}
\textbf{Gợi ý để dạy và tự nghĩ / Teaching and reflection prompts}

1. Vì sao nhiều trẻ chỉ được mềm lại khi đã quá mệt? \textit{Why are many children treated softly only when they are already exhausted?}

2. Điều đó dạy trẻ bài học gì? \textit{What lesson does that teach a child?}

3. Mình có thể dịu hơn từ lúc nào? \textit{How much earlier can we become gentle?}
\end{tuhocnhanh}
\ngancach

\section{Em Ngoan Có Phải Vì Em Được Yêu Không}
\begin{truyen}{Em Ngoan Có Phải Vì Em Được Yêu Không}{Am I Well-Behaved Because I Am Loved}
\chuhoa{C}{ó những học sinh ngoan quá mức. Không dám làm phiền, không dám xin thêm, không dám nói mình buồn.}
\textit{(Some students are good to an extreme. They do not dare to bother anyone, ask for more, or say they are sad.)}

Bề ngoài người lớn rất thích kiểu ngoan này.
\textit{(Adults often like this kind of obedience.)}

Nhưng có khi đằng sau là một câu hỏi rất sâu: nếu em không dễ chịu nữa, người lớn còn giữ em không?
\textit{(But behind it may be a deep question: if I stop being easy, will adults still keep me?) }

Ngoan đôi khi là cách một đứa trẻ tự bảo vệ chỗ đứng của mình trong tình thương.
\textit{(Good behavior can sometimes be a child's way of protecting their place in love.)}

\end{truyen}

\begin{baihoc}
Không phải sự ngoan nào cũng đến từ bình an. Có sự ngoan đi ra từ nỗi sợ bị mất chỗ.
\textit{(Not all obedience comes from peace. Some comes from fear of losing one's place.)}
\end{baihoc}

\begin{ghinhoanh}
\textbf{Short English to remember:}

Not all obedience comes from peace.

\end{ghinhoanh}

\begin{tuhocnhanh}
\textbf{Gợi ý để dạy và tự nghĩ / Teaching and reflection prompts}

1. Vì sao một số học sinh ngoan quá mức? \textit{Why are some students obedient to an extreme?}

2. Câu hỏi chìm dưới sự ngoan này là gì? \textit{What hidden question lies under this obedience?}

3. Mình có đang tạo ra cảm giác phải ngoan mới được yêu không? \textit{Are we creating the feeling that love must be earned by behavior?}
\end{tuhocnhanh}
\ngancach

\section{Lúc Em Tệ Nhất, Ai Ở Lại}
\begin{truyen}{Lúc Em Tệ Nhất, Ai Ở Lại}{Who Stays When I Am at My Worst}
\chuhoa{M}{ột đứa trẻ thường nhớ rất lâu người đã ở lại với mình đúng lúc mình xấu xí nhất.}
\textit{(A child remembers for a long time the person who stayed at the moment they looked worst.)}

Có em bị bắt gặp nói dối, làm sai, hoặc cư xử dở.
\textit{(Some are caught lying, doing wrong, or behaving badly.)}

Nếu tất cả người lớn chỉ dùng lúc đó để chốt tội, đứa trẻ sẽ học rằng sai lầm làm mình mất hết chỗ dựa.
\textit{(If every adult uses that moment only to condemn, the child learns that mistakes remove all support.)}

Nhưng nếu có ai vừa sửa vừa ở lại, bài học sẽ sâu hơn và lành hơn.
\textit{(But if someone both corrects and stays, the lesson becomes deeper and healthier.)}

\end{truyen}

\begin{baihoc}
Trẻ con rất cần người không bỏ mình đúng lúc mình đáng bị bỏ nhất.
\textit{(Children deeply need someone who does not leave when they seem most leaveable.)}
\end{baihoc}

\begin{ghinhoanh}
\textbf{Short English to remember:}

Children remember who stays when they seem least lovable.

\end{ghinhoanh}

\begin{tuhocnhanh}
\textbf{Gợi ý để dạy và tự nghĩ / Teaching and reflection prompts}

1. Vì sao lúc sai là lúc trẻ cần người ở lại nhất? \textit{Why is a child's worst moment often the time they most need someone to stay?}

2. Sửa và bỏ đi khác sửa và ở lại ở đâu? \textit{How is correction with abandonment different from correction with presence?}

3. Mình từng là kiểu người lớn nào trong những lúc như vậy? \textit{What kind of adult have I been in those moments?}
\end{tuhocnhanh}
\ngancach

\section{Con Điểm Này Thì Mẹ Còn Thương Con Không}
\begin{truyen}{Con Điểm Này Thì Mẹ Còn Thương Con Không}{Will Mom Still Love Me with This Score}
\chuhoa{C}{ó em cầm bài kiểm tra, đứng ngoài cửa lớp, hỏi rất nhỏ: “Con điểm này thì mẹ còn thương con không thầy?”}
\textit{(Some students hold a test paper at the classroom door and ask quietly, “With this score, will my mom still love me?”)}

Nghe xong mới thấy nhiều đứa trẻ không sợ điểm trước tiên.
\textit{(In that moment we see that many children do not fear the score first.)}

Các em sợ cách người lớn đổi mặt.
\textit{(They fear the change in an adult's face.)}

Một đứa trẻ hỏi như vậy là đã sống quá gần cảm giác tình thương có điều kiện.
\textit{(A child asking that has already lived too close to conditional love.)}

\end{truyen}

\begin{baihoc}
Khi trẻ hỏi mình còn được thương không, vấn đề lớn hơn bài kiểm tra rất nhiều.
\textit{(When a child asks whether they are still loved, the problem is far bigger than the test.)}
\end{baihoc}

\begin{ghinhoanh}
\textbf{Short English to remember:}

The real fear is often not the grade, but the loss of love.

\end{ghinhoanh}

\begin{tuhocnhanh}
\textbf{Gợi ý để dạy và tự nghĩ / Teaching and reflection prompts}

1. Câu hỏi này cho thấy đứa trẻ đang sống trong cảm giác gì? \textit{What feeling does this question reveal?}

2. Vì sao đổi sắc mặt của người lớn lại đáng sợ đến vậy? \textit{Why is a change in an adult's face so frightening?}

3. Người dạy nên đáp câu này bằng giọng nào? \textit{How should a teacher answer this question?}
\end{tuhocnhanh}
\ngancach

\section{Em Ngoan Vì Em Sợ Bị Chán}
\begin{truyen}{Em Ngoan Vì Em Sợ Bị Chán}{I Behave Because I Fear Being Given Up On}
\chuhoa{C}{ó em lúc nào cũng “dạ”, cũng “vâng”, cũng cười cho vừa ý người lớn.}
\textit{(Some students are always saying “yes”, always smiling in the way adults prefer.)}

Tôi hỏi: “Sao em không bao giờ cãi lại?”
\textit{(I asked, “Why do you never argue back?”)}

Em đáp rất nhanh: “Tại em sợ người lớn chán em.”
\textit{(The student answered quickly, “Because I'm afraid adults will get tired of me.”)}

Một câu như vậy làm sự ngoan bỗng nghe không còn nhẹ nữa.
\textit{(A line like that makes obedience stop sounding harmless.)}

\end{truyen}

\begin{baihoc}
Có những đứa trẻ ngoan không phải vì yên lòng, mà vì đang lo mất chỗ trong mắt người lớn.
\textit{(Some children behave well not from peace, but from fear of losing their place in adult eyes.)}
\end{baihoc}

\begin{ghinhoanh}
\textbf{Short English to remember:}

Some children behave well because they are afraid of being given up on.

\end{ghinhoanh}

\begin{tuhocnhanh}
\textbf{Gợi ý để dạy và tự nghĩ / Teaching and reflection prompts}

1. Vì sao sự ngoan này nghe buồn hơn nghe mừng? \textit{Why does this kind of obedience sound sadder than comforting?}

2. Trẻ đang giữ điều gì bằng cách luôn làm vừa lòng? \textit{What is the child trying to keep by always pleasing adults?}

3. Mình có đang thưởng quá mạnh cho kiểu ngoan do sợ hãi không? \textit{Am I rewarding fear-based obedience too strongly?}
\end{tuhocnhanh}
\ngancach

\section{Lúc Em Khóc Người Lớn Mới Tin}
\begin{truyen}{Lúc Em Khóc Người Lớn Mới Tin}{Adults Believe Me Only When I Cry}
\chuhoa{C}{ó em nói mệt ba lần không ai để ý. Đến lần thứ tư em khóc, người lớn mới dịu xuống.}
\textit{(Some students say they are tired three times and no one notices. On the fourth time they cry, adults finally soften.)}

Trẻ con học rất nhanh kiểu quy luật này.
\textit{(Children learn this rule very quickly.)}

Các em hiểu rằng muốn được tin, phải đau ra mặt.
\textit{(They understand that to be believed, pain must become visible.)}

Đó là cách rất mệt để lớn lên.
\textit{(That is an exhausting way to grow up.)}

\end{truyen}

\begin{baihoc}
Nếu người lớn chỉ tin khi trẻ đã vỡ ra, trẻ sẽ học cách chịu quá lâu trước khi cầu cứu.
\textit{(If adults believe children only after they break down, children learn to endure far too long before asking for help.)}
\end{baihoc}

\begin{ghinhoanh}
\textbf{Short English to remember:}

Children should not have to break down before being believed.

\end{ghinhoanh}

\begin{tuhocnhanh}
\textbf{Gợi ý để dạy và tự nghĩ / Teaching and reflection prompts}

1. Vì sao trẻ hay chịu quá lâu rồi mới vỡ? \textit{Why do children endure too long before breaking?}

2. Người lớn đang bỏ sót những tín hiệu nào trước khi trẻ khóc? \textit{What signals are adults missing before tears appear?}

3. Mình có đang đợi quá muộn mới tin học sinh không? \textit{Am I believing students too late?}
\end{tuhocnhanh}
\ngancach

\section{Em Có Được Giận Không}
\begin{truyen}{Em Có Được Giận Không}{Am I Allowed to Be Angry}
\chuhoa{C}{ó học sinh hỏi rất thật: “Nếu em giận thì có phải là em hư không?”}
\textit{(Some students ask very honestly, “If I get angry, does that make me bad?”)}

Nhiều đứa trẻ chỉ được phép vui, ngoan, và biết ơn.
\textit{(Many children are only allowed to be cheerful, obedient, and grateful.)}

Còn giận, bực, khó chịu thì bị coi là hỗn.
\textit{(Anger, frustration, and irritation are treated as disrespect.)}

Khi một cảm xúc tự nhiên cũng bị cấm, trẻ sẽ không biết sống thật với mình bằng cách nào.
\textit{(When a natural feeling is forbidden, children no longer know how to live honestly with themselves.)}

\end{truyen}

\begin{baihoc}
Yêu thương đúng không phải là cấm trẻ có cảm xúc khó, mà là dạy trẻ đi qua nó cho lành.
\textit{(Healthy love does not forbid difficult feelings. It teaches children how to move through them safely.)}
\end{baihoc}

\begin{ghinhoanh}
\textbf{Short English to remember:}

Healthy love does not forbid difficult feelings.

\end{ghinhoanh}

\begin{tuhocnhanh}
\textbf{Gợi ý để dạy và tự nghĩ / Teaching and reflection prompts}

1. Vì sao nhiều trẻ sợ cả cảm xúc giận của mình? \textit{Why are many children afraid of their own anger?}

2. Cấm cảm xúc khác với dạy cảm xúc ở đâu? \textit{How is forbidding feelings different from teaching them?}

3. Mình có đang chỉ cho học sinh được phép có vài cảm xúc đẹp không? \textit{Am I allowing students only the feelings that look nice?}
\end{tuhocnhanh}
\ngancach

\section{Sau Khi Mắng, Người Lớn Có Ôm Lại Em Không}
\begin{truyen}{Sau Khi Mắng, Người Lớn Có Ôm Lại Em Không}{After Scolding, Do Adults Draw Close Again}
\chuhoa{C}{ó em kể: “Lúc mẹ la thì lớn lắm. Nhưng la xong mẹ im luôn.”}
\textit{(One student said, “When Mom scolds me, it is loud. But after that she stays silent.”)}

Nhiều đứa trẻ không sợ nhất ở cơn giận.
\textit{(Many children do not fear the outburst most.)}

Các em sợ khoảng lạnh kéo dài sau cơn giận.
\textit{(They fear the cold distance that comes after it.)}

Nếu người lớn không quay lại, trẻ sẽ phải tự ôm lấy phần hoảng của mình.
\textit{(If adults do not return, children must hold their own fear alone.)}

\end{truyen}

\begin{baihoc}
Quay lại sau khi nóng là một phần rất thật của tình thương có trách nhiệm.
\textit{(Coming back after anger is a very real part of responsible love.)}
\end{baihoc}

\begin{ghinhoanh}
\textbf{Short English to remember:}

Returning after anger is part of responsible love.

\end{ghinhoanh}

\begin{tuhocnhanh}
\textbf{Gợi ý để dạy và tự nghĩ / Teaching and reflection prompts}

1. Vì sao khoảng lạnh sau cơn nóng làm trẻ sợ? \textit{Why does the cold silence after anger frighten children?}

2. Quay lại sau khi mắng có ý nghĩa gì? \textit{What does it mean to return after scolding?}

3. Mình có biết nối lại với học sinh sau lúc căng không? \textit{Do I know how to reconnect after tension?}
\end{tuhocnhanh}
\ngancach

\section{Một Câu Thôi, Nhưng Cứu Được Cả Buổi Tối}
\begin{truyen}{Một Câu Thôi, Nhưng Cứu Được Cả Buổi Tối}{One Sentence Can Save the Whole Evening}
\chuhoa{C}{ó lần tôi chỉ nói với một em đang run tay cầm bài kiểm tra: “Không sao, về nhà rồi mình tính tiếp.”}
\textit{(Once I said to a student whose hands were shaking over a test paper, “It is okay. We'll figure it out after you go home.”)}

Em thở ra rất rõ.
\textit{(The student visibly exhaled.)}

Có những đứa trẻ không cần bài giảng dài lúc đó.
\textit{(Some children do not need a long lecture in that moment.)}

Các em chỉ cần biết tối nay mình chưa bị bỏ rơi.
\textit{(They only need to know they are not being abandoned tonight.)}

\end{truyen}

\begin{baihoc}
Đôi khi tình thương đi vào đời trẻ bằng một câu ngắn nhưng đúng lúc.
\textit{(Sometimes love enters a child's life through one short sentence said at the right time.)}
\end{baihoc}

\begin{ghinhoanh}
\textbf{Short English to remember:}

One timely sentence can calm a frightened child.

\end{ghinhoanh}

\begin{tuhocnhanh}
\textbf{Gợi ý để dạy và tự nghĩ / Teaching and reflection prompts}

1. Vì sao một câu ngắn đúng lúc lại có sức lớn? \textit{Why can one short sentence carry such weight?}

2. Học sinh đang cần nghe gì nhất trong khoảnh khắc sợ hãi? \textit{What do students most need to hear in a fearful moment?}

3. Mình có đang nói quá nhiều thay vì nói đúng không? \textit{Am I speaking too much instead of speaking rightly?}
\end{tuhocnhanh}
